\documentclass[11pt]{article}

\usepackage[a4paper,margin=1in]{geometry}
\usepackage{amsmath,amssymb,amsthm}
\usepackage{microtype}

\newtheorem{theorem}{Theorem}
\newtheorem{lemma}[theorem]{Lemma}
\newtheorem{remark}[theorem]{Remark}

\newcommand{\dd}{\,\mathrm d}

\title{A Partial Cone Lift from the House Graph to Atlas 196}
\author{}
\date{}

\begin{document}

\maketitle

\section{Statement}

Let $F$ be the house graph, Atlas $43$, and let
\[
 H=K_1\vee F.
\]
The graph $H$ is Atlas $196$, with graph6 code \texttt{ER\string~g}.  Put
\[
 \phi(z):=\Phi_F(z)
 =z(2z-1)(3z^2-3z+1)
 =6z^4-9z^3+5z^2-z.
\]
The improved global house theorem in
\texttt{notes/atlas43\_improved\_global\_interval.tex} proves
\begin{equation}\label{eq:house-input}
 t(F,V)\geq\phi(t(K_2,V))
 \qquad\text{whenever}\qquad
 t(K_2,V)\geq
 a:=0.7085121949942608482\ldots,
\end{equation}
where (a) is the unique root in ((3/4,1)) of
\[
 5440a^4-10848a^3+7981a^2-2580a+309=0.
\]

For $c\in[a,1]$, define the value at $a$ of the tangent to $\phi$ at $c$:
\[
 L(c):=\phi(c)+(a-c)\phi'(c).
\]
There is a unique $c_0\in(a,1)$ satisfying
\begin{equation}\label{eq:c0}
 L(c_0)=0.
\end{equation}
Define
\begin{equation}\label{eq:p0}
 p_0:=\frac1{2-c_0}.
\end{equation}

\begin{theorem}[Atlas 196 partial range]\label{thm:atlas196}
Every graphon $W$ of edge density $p\geq p_0$ satisfies
\[
 t(H,W)\geq\Phi_H(p).
\]
Numerically,
\[
\begin{split}
 c_0&=0.8412137423978822016\ldots,\\
 p_0&=0.8629719186257049718\ldots,
\end{split}
\]
and in particular the theorem holds throughout the certified rounded range
\[
 \boxed{p\geq0.8629719187.}
\]
\end{theorem}

\section{A partial-range lifting lemma}

The usual universal-vertex lifting theorem assumes that the base inequality
starts at a zero of its target.  Only one line of its proof needs to be
changed when the known base range begins later.

\subsection{The weighted link estimate}

We include the weighted local Goodman inequality used in the lift so that
the argument is self-contained.  Write
\[
 d(x)=\int W(x,y)\dd y,
 \qquad
 \tau(x)=\int W(x,y)W(x,z)W(y,z)\dd y\dd z.
\]

\begin{lemma}[Weighted rooted-triangle inequality]\label{lem:weighted-goodman}
If \(W\) has edge density \(p>0\), then, for every integer \(k\geq0\),
\[
 \int d(x)^k\tau(x)\dd x
 \geq
 \frac{2p-1}{p}\int d(x)^{k+2}\dd x.
\]
Consequently, for every integer \(n\geq2\),
\begin{equation}\label{eq:weighted-link}
 \int d(x)^n z_x\dd x
 \geq
 \left(2-\frac1p\right)\int d(x)^n\dd x,
\end{equation}
where \(z_x=\tau(x)/d(x)^2\) when \(d(x)>0\), and \(z_x=0\)
when \(d(x)=0\).
\end{lemma}

\begin{proof}
Put \(M_j=\int d^j\), and set
\[
 I_k=\int d(x)^k\tau(x)\dd x,
 \qquad
 J_k=\int d(x)^kW(x,y)d(y)\dd x\dd y.
\]
The pointwise inequality \(ab\geq a+b-1\), applied to
\(W(x,y)W(x,z)\), gives
\begin{equation}\label{eq:I-J}
 I_k\geq2J_k-pM_k.
\end{equation}

Let \(U=1-W\), \(u(x)=\int U(x,y)\dd y=1-d(x)\), and
\(q=1-p\).  If \(T_U\) is the integral operator with kernel \(U\), then
\[
 T_Wd=d-q+T_Uu,
\]
so
\[
 J_k=M_{k+1}-qM_k+\int d(x)^k(T_Uu)(x)\dd x.
\]
Symmetry of \(U\) gives
\[
\begin{split}
 &2\left[
 \int d(x)^k(T_Uu)(x)\dd x-\int d(x)^ku(x)^2\dd x
 \right]\\
 &\qquad=
 \int U(x,y)\bigl(u(y)-u(x)\bigr)
 \bigl(d(x)^k-d(y)^k\bigr)\dd x\dd y
 \geq0.
\end{split}
\]
The last sign holds because \(u=1-d\), so the two displayed differences
have the same sign.  Therefore
\[
 J_k\geq
 M_{k+1}-qM_k+\int d^k(1-d)^2
 =M_{k+2}-M_{k+1}+pM_k.
\]
Combining this with \eqref{eq:I-J} yields
\[
 I_k\geq2M_{k+2}-2M_{k+1}+pM_k.
\]
Finally,
\[
\begin{split}
 &2M_{k+2}-2M_{k+1}+pM_k
 -\frac{2p-1}{p}M_{k+2}\\
 &\qquad=
 \frac1p\int d(x)^k\bigl(d(x)-p\bigr)^2\dd x
 \geq0.
\end{split}
\]
This proves the first assertion.  Since
\(d(x)^nz_x=d(x)^{n-2}\tau(x)\), applying it with \(k=n-2\)
proves \eqref{eq:weighted-link}.
\end{proof}

\subsection{The tangent lift}

\begin{lemma}[Tangent lift from a partial base range]\label{lem:lift}
Let $G$ be a graph on $n\geq2$ vertices and let
$\psi(z)=\Phi_G(z)$.  Suppose that, for some $a\in[0,1)$,
\begin{enumerate}
 \item every graphon $V$ of density $z\geq a$ satisfies
       $t(G,V)\geq\psi(z)$;
 \item $\psi,\psi'$, and $\psi''$ are nonnegative on $[a,1]$.
\end{enumerate}
Let $W$ have density $p<1$, set
\[
 c=2-\frac1p,
\]
and suppose
\[
 c\in[a,1],
 \qquad
 \psi(c)+(a-c)\psi'(c)\leq0.
\]
Then
\[
 t(K_1\vee G,W)\geq\Phi_{K_1\vee G}(p).
\]
\end{lemma}

\begin{proof}
Write $d(x)=\int W(x,y)\dd y$.  For $d(x)>0$, let $W_x$ be the link
graphon on the probability measure $W(x,y)\dd y/d(x)$, and let $z_x$ be
its edge density; when $d(x)=0$, set $z_x=0$.  The link identity gives
\begin{equation}\label{eq:link}
 t(K_1\vee G,W)
 =\int d(x)^n t(G,W_x)\dd x.
\end{equation}

Let
\[
 \ell_c(z)=\psi(c)+\psi'(c)(z-c)
\]
be the tangent at $c$.  If $z\geq a$, convexity and the base theorem give
\[
 t(G,V)\geq\psi(z)\geq\ell_c(z).
\]
If $z<a$, then $\psi'(c)\geq0$ and the assumed tangent condition give
\[
 \ell_c(z)\leq\ell_c(a)\leq0\leq t(G,V).
\]
Consequently $t(G,V)\geq\ell_c(z)$ for every graphon $V$, irrespective of
whether its density lies in the proved base range.

Substitute this affine minorant into \eqref{eq:link}.
Lemma~\ref{lem:weighted-goodman} gives
\[
 \int d(x)^n z_x\dd x
 \geq
 \left(2-\frac1p\right)\int d(x)^n\dd x
 =c\int d(x)^n\dd x.
\]
Since $\psi'(c)\geq0$, the tangent term is nonnegative after integration,
and hence
\[
 t(K_1\vee G,W)
 \geq
 \psi(c)\int d(x)^n\dd x.
\]
Jensen gives $\int d^n\geq p^n$, and $\psi(c)\geq0$, so
\[
 t(K_1\vee G,W)\geq p^n\psi(c).
\]
Finally, writing $q=1-p$, one has
\[
 1-c=\frac qp,
 \qquad
 p^n\psi(c)=q^n\chi_G\!\left(\frac pq\right)
 =\Phi_{K_1\vee G}(p),
\]
where the last identity uses
$\chi_{K_1\vee G}(x)=x\chi_G(x-1)$.  This proves the lemma.  The endpoint
$p=1$ follows directly from $W=1$ almost everywhere.
\end{proof}

\section{Verification for the house target}

Differentiation gives
\[
 \phi'(z)=24z^3-27z^2+10z-1,
 \qquad
 \phi''(z)=2(3z-1)(12z-5).
\]
Thus $\phi''(z)>0$ for $z\geq1/2$.  Also
$\phi'(1/2)=1/4$, so $\phi'$ is positive on $[1/2,1]$, and
$\phi$ is nonnegative there because its only vanishing factor in that
interval is $2z-1$.  The analytic hypotheses of
Lemma~\ref{lem:lift} therefore hold on $[a,1]$.

Moreover,
\[
 L'(c)=(a-c)\phi''(c)<0
 \qquad(a<c<1).
\]
At the two endpoints,
\[
 L(a)=\phi(a)>0,
 \qquad
 L(1)=1+6(a-1)=6a-5<0.
\]
Hence the zero $c_0$ in \eqref{eq:c0} exists and is unique, and
$L(c)\leq0$ for every $c\geq c_0$.

The map $p\mapsto2-1/p$ is increasing.  Therefore $p\geq p_0$ implies
$c\geq c_0$, and Lemma~\ref{lem:lift}, with the accepted house input
\eqref{eq:house-input}, proves Theorem~\ref{thm:atlas196}.

The numerical values are obtained by first isolating the unique quartic
root (a), then using the strict monotonicity of (L) to isolate the
unique zero in \eqref{eq:c0}.  Exact rational interval arithmetic gives
\[
 0.8629719186<p_0<0.8629719187,
\]
so the rounded decimal threshold in the theorem is on the proved side.

\section{Scope}

This note uses the improved global house theorem only on its proved range.  It
does not assume the unresolved house inequality below $a$, and it applies
to arbitrary graphons.  It leaves
\[
 \frac23\leq p<p_0
\]
open for Atlas 196, so the catalogue status is partial rather than
positive.  A future full proof for the house would satisfy the ordinary
lifting theorem with $a=1/2$ and would immediately settle Atlas 196 on its
entire required interval.

\end{document}
